% Modernised reading — fonds Alexandre Grothendieck, shelfmark 59, pages 1-27.
%
% This is an INTERPRETATION, not a transcription. Everything the critical
% apparatus carried moves into footnotes. In any dispute of substance the
% transcription governs: whoever cites Grothendieck must cite batch-01.fr.tex
% and batch-02.fr.tex.
%
% Pass: Opus 5.5 (claude-opus-5-5), 2026-10-04 — first pass, unchecked against
% the pages by a human. The folder was read whole: two batches, pages 1-20 and
% 21-27; page 21 is blank. Two runs: pp. 1-20 (paragraphs a) to f), with e)
% and f) each used twice) and pp. 22-27 (his own pagination 1-6).
%
% Notation fixed for the whole document. Pic, Alb, NS without his underlines
% (they are representable functors throughout). B = Pic^0_{X/S}, A =
% Pic^0_{B/S} = Alb^0_{X/S}. For a class lambda in NS, lambda~ is the
% homomorphism with Mumford's convention lambda~(x) = t_x^*M (x) M^{-1}. For a
% B-torsor T and u : B -> A, u_*T = T x^B A; T^{(m)} = [m]_*T, T^{(-1)} the
% inverse torsor. The Weil sheaf of the pages is called the Poincaré sheaf.
% The Theta-isomorphism of p. 8 (his c) is written phi_Theta, to keep c for
% the Abel-Jacobi isomorphism of p. 23. delta' = -phi(delta) throughout
% (p. 10 once writes delta' = phi(delta)).
%
% Substantive departures from the pages, each footnoted:
% — p. 3, the twisted Picard torsor is written with the INVERSE of the
%   pushed torsor (the page has none); with Mumford's convention, which is the
%   one of (**) p. 18, this is what holds. It makes (*) p. 11 come out exactly
%   as the page boxes it, and agrees with Mukai's transform.
% — p. 4, the arrow from Alb^{-chi} and the box in Alb^{chi} are not in
%   conflict (the transcription flags them): a torsor map from Alb^{-chi} is a
%   section of Alb^{chi} x T.
% — p. 8, phi(delta) for a curve is MINUS the canonical polarisation (Mattuck;
%   checked on an elliptic curve); the page writes plus. The product torsor
%   P^{2g-2} of p. 8 is not confirmed: our count gives exponent 0; flagged as
%   a doubt, not asserted.
% — p. 10, chi(delta') = chi(delta)^{n-1} from degrees is only up to sign;
%   supplied by density of the ample cone. phi(phi(delta)) carries (-1)^n.
%   The converse « phi(delta) non dégénérée => chi(delta) != 0 » needs n >= 2.
% — p. 10 N.B.: the expected indivisibility of delta' fails for types such as
%   (2,2); footnoted.
% — p. 15, Todd(curve) = 1 - K/2 (the page writes K/2, degree 1-g); general
%   rank rho restored, the page has rho = 1.
% — p. 16, ell' is not additive in Pic (the page's parenthesis); what holds is
%   (**). p. 19's phi(delta)~(-beta) contradicts (**) unless the convention
%   for lambda~ flips; (**) kept.
% — p. 26 c): Theta' for the page's Theta.
%
% What the folder announces and never establishes: the expected value K of
% p. 8 (« à prouver ! ») and the N.B. on the moduli of curves (strong
% Franchetta); the N.B. of p. 10 on moduli of polarised abelian schemes; the
% compatibilities of trivialisations asked for in the margins of pp. 11-12
% and 27, on which the conclusions of pp. 12-13 depend; the universal
% property of Alb^1 (« on prouve », p. 22); the question of p. 23, left
% open after p. 24's two lines; the study of « la petite topologie de
% Albanese » (p. 20, reading doubtful).
\documentclass[11pt,a4paper]{article}
\input{../preamble/grothendieck}

\folder{59}
\pages{1}{27}
\dating{[années 1960-1970]}
\watermark{Édition de démonstration\\interprétation personnelle de l'œuvre}
\foldertitle{Construction des faisceaux inversibles sur schémas de Picard : notes manuscrites (s.d.) — lecture modernisée du dossier entier}

\begin{document}

\section*{Résumé}

\begin{resume}
Une courbe, ou plus généralement une variété, et l'espace qui paramètre ses
fibrés en droites : comment fabriquer, sans rien choisir, des fibrés en droites
sur cet espace-là ?

Sur une variété $X$ vivent des fibrés en droites, et ceux-ci forment à leur
tour un espace géométrique, le \emph{schéma de Picard}. Sa composante neutre,
quand tout va bien, est une \emph{variété abélienne} — un objet dont les points
s'additionnent, l'analogue algébrique d'un tore — et cette variété abélienne
possède une \emph{duale}, qui paramètre à son tour ses propres fibrés en
droites. On a donc une tour : $X$, puis son Picard, puis le Picard de son
Picard. Le dossier cherche à descendre cette tour d'un étage de façon
canonique, c'est-à-dire sans choisir de point base, de diviseur ni de
coordonnées. Le problème vient de ce que presque toutes les constructions
naturelles demandent un choix — un point de $X$ pour normaliser, un fibré de
référence — et que le résultat change avec le choix.

L'idée est la plus simple qui soit : à chaque fibré $L$ sur $X$ on associe sa
cohomologie, et de la cohomologie on ne garde qu'un nombre — en fait une
droite, son \emph{déterminant}. Quand $L$ varie, ces droites forment un fibré
en droites sur le schéma de Picard. Il dépend encore du point base choisi, mais
d'une façon que l'on contrôle exactement : le changement est une puissance,
d'exposant la caractéristique d'Euler, d'un fibré que l'on sait décrire. Toute
la suite du dossier exploite cette dépendance maîtrisée. On en tire une
application entre les « classes » de fibrés de $X$ et celles de la variété
duale — pas un homomorphisme, un polynôme — et, au bout d'une comptabilité de
torseurs (des espaces qui ressemblent à un groupe, mais sans origine
privilégiée), une \emph{section canonique} d'un certain torseur : un objet que
l'on n'avait pas le droit d'espérer sans choix, et qui est là.

Deux exemples éclairent la construction. Pour une courbe, l'application
retrouve, au signe près, la polarisation classique de la jacobienne — le
diviseur thêta — et la section canonique attendue est le fibré canonique de la
courbe ; le dossier le conjecture, et ses dernières pages l'obtiennent par un
autre chemin, à partir de la symétrie du diviseur thêta. Pour une variété
abélienne, le calcul passe par ce qu'on appelle aujourd'hui la
\emph{transformation de Fourier} des cycles, et produit une « inversion » des
polarisations : chaque polarisation de la variété en donne une sur la duale,
et la composée des deux est la multiplication par la caractéristique d'Euler.
Il en sort des énoncés de trivialité canonique pour certaines puissances de
torseurs. Une dernière construction recommence avec un complexe de faisceaux
quelconque à la place d'un fibré, et relit tout comme le déterminant d'une
transformation de type Fourier.

Les feuillets sont rapides, et presque chaque page dit son inquiétude sur les
signes : « vérifier si le signe est correct ! », « attention aux signes ?? »,
« question de signe ! ». L'inquiétude était fondée — plusieurs signes sont à
retourner, et l'un des résultats intermédiaires ne se confirme pas en l'état ;
les notes de bas de page disent lesquels. Ce qui reste juste, et c'est
l'essentiel, est la forme de la construction : une seule idée, le déterminant
de la cohomologie, suivie jusqu'au bout et vérifiée sur les deux exemples où
l'on sait calculer.

\keywords{determinant of cohomology, Picard scheme, Néron-Severi group scheme, Poincaré bundle, torsor under an abelian scheme, contracted product, Albanese torsor, Grothendieck-Riemann-Roch, Picard bundle, Fourier-Mukai transform, Fourier transform on Chow groups, dual polarisation, principally polarised abelian scheme, theta divisor, Riemann singularity theorem, autoduality of the Jacobian, canonical bundle, Franchetta conjecture}
\end{resume}

\pagerange{1}{27}
\section*{Le dossier, sa charpente et ses notations}

Vingt-sept pages, dont une blanche (la page 21). Elles forment deux suites
distinctes. Les pages 1 à 20 sont une rédaction continue, en paragraphes
lettrés a) à f) — avec deux e) (pages 5 et 14) et deux f) (pages 9 et 16) ;
on les distingue ici par leur contenu et leurs pages. Les pages 22 à 27,
numérotées de 1 à 6 de sa main, sont une reprise qui commence par fixer une
convention de signe dont la première suite avait besoin, puis passe aux
schémas abéliens principalement polarisés.

L'argument se lit en huit stations :

\begin{itemize}
\item le fibré déterminant $M_g$ sur le schéma de Picard, et sa dépendance en
  la section $g$ (pages 1 et 2) ;
\item l'application $\varphi : \mathrm{NS}_{X/S} \to \mathrm{NS}_{B/S}$, la
  torsion de $\mathrm{Pic}_{B/S}$ par $P^{\delta}$ et la section canonique
  $L(\delta)$ (pages 2 à 5) ;
\item le calcul de $\varphi(\delta)$ par Riemann--Roch (pages 5 à 7) ;
\item premier exemple, les courbes (pages 7 à 9) ;
\item second exemple, les schémas abéliens : transformation de Fourier,
  inversion des polarisations, l'isomorphisme $(*)$ et les trivialisations de
  puissances de torseurs (pages 9 à 13) ;
\item la généralisation par un complexe parfait $\mathcal{F}$ (pages 14 à 16) ;
\item une seconde interprétation, par le déterminant d'une transformation de
  Fourier (pages 16 à 20) ;
\item la reprise : le torseur d'Albanese et sa convention de signe (pages 22 à
  24), puis le torseur thêta d'un schéma abélien principalement polarisé et sa
  section $K$ (pages 25 à 27).
\end{itemize}

\emph{Notations, valables pour tout le dossier.} $f : X \to S$ est propre,
plat, de présentation finie, avec $f_{*}\mathcal{O}_X = \mathcal{O}_S$
universellement ; $P = \mathrm{Pic}_{X/S}$, et l'on suppose que
$B = P^{0} = \mathrm{Pic}^{0}_{X/S}$ est un schéma abélien. On note
$A = \mathrm{Pic}^{0}_{B/S}$ son dual, que la page 3 identifie à
$\mathrm{Alb}^{0}_{X/S}$. Pour une section $\delta$ de
$\mathrm{NS}_{X/S}$, $P^{\delta}$ est la composante correspondante de $P$,
torseur sous $B$. Les foncteurs $\mathrm{Pic}$, $\mathrm{Alb}$, $\mathrm{NS}$
sont soulignés sur les pages, comme partout dans le fonds ; on ne les souligne
pas ici. Le « faisceau de Weil » des pages est le \emph{faisceau de Poincaré}
normalisé le long d'une section.

Pour une classe $\lambda$ de $\mathrm{NS}$ d'un schéma abélien, on note
$\tilde{\lambda}$ l'homomorphisme vers le dual avec la convention de Mumford :
$\tilde{\lambda}(x) = [\,t_x^{*}M \otimes M^{-1}\,]$ pour $M$ de classe
$\lambda$.\footnote{Les pages ne fixent jamais cette convention, et c'est
d'elle que dépend la moitié des signes qu'elles signalent comme douteux. On
prend celle de Mumford parce que c'est celle sous laquelle la formule $(**)$ de
la page 18 est juste telle que la page l'écrit.} Pour un $B$-torseur $T$ et un
homomorphisme $u : B \to A$, on écrit $u_{*}T = T \times^{B} A$ le
$A$-torseur poussé — les pages écrivent $T \overset{B}{\times} (A, u)$ ;
$T \times^{A} T'$ est le produit contracté de deux $A$-torseurs ;
$T^{(m)} = [m]_{*}T$, et $T^{(-1)}$ est le torseur inverse. Enfin
$\mathrm{Alb}^{k}_{X/S} = (\mathrm{Alb}^{1}_{X/S})^{(k)}$.

Deux homonymies sont à défaire dès maintenant. La page 8 appelle $c$
l'isomorphisme $B \to A$ défini par le diviseur $\Theta$ d'une jacobienne ; la
page 23 appelle $c$ l'isomorphisme d'Abel--Jacobi. Ce ne sont pas les mêmes —
la page 23 pense qu'ils diffèrent d'un signe, et c'est exact — et l'on écrit
$\phi_{\Theta}$ pour le premier, $c$ pour le second. D'autre part la lettre
$\delta$ désigne dans les pages 1 à 20 une classe de Néron--Severi, et à la
page 22 l'application $X \times_S X \to A$, $(x, y) \mapsto \psi(y) -
\psi(x)$ ; on écrit $\Delta_{\psi}$ pour celle-ci.

\pagerange{1}{2}
\section*{Le fibré déterminant et le changement de section (pages 1 et 2)}

Supposons d'abord que $X/S$ ait une section $g$. Sur $X \times_S P$ vit alors
le faisceau de Poincaré $\mathcal{L}_g$, normalisé, c'est-à-dire trivialisé le
long de $g(S) \times_S P$. Son image directe dérivée par
$f_P : X \times_S P \to P$ est un complexe parfait, et l'on pose
\[
M_g = \det\, R f_{P*}(\mathcal{L}_g),
\]
module inversible sur $P$ : c'est le \emph{fibré déterminant de la
cohomologie}.\footnote{La page écrit ${\det}^{*}$, et ajoute après
« module inversible » un mot lu « gradué » sans certitude, suivi d'un mot
illisible. Le dossier se sert du déterminant d'un complexe parfait sans le
construire ; la construction publiée en est celle de Knudsen et Mumford (1976).
La datation du dossier, « années 1960-1970 », ne permet pas de dire laquelle
des deux est antérieure, et rien ici ne le décide.}

Changeons de section. Le faisceau normalisé le long de $g'$ est
\[
\mathcal{L}_{g'} \simeq \mathcal{L}_g \otimes f_P^{*}\bigl(g'^{*}_P \mathcal{L}_g\bigr)^{-1},
\]
d'où, par la formule de projection et la multiplicativité du déterminant, un
isomorphisme canonique
\[
(*_{b}) \qquad M_{g'} \simeq M_g \otimes \bigl(g'^{*}_P \mathcal{L}_g\bigr)^{-\chi}
= M_g \otimes \bigl(N_g N_{g'}^{-1}\bigr)^{\chi},
\]
où $\chi : P \to \mathbb{Z}$ est la fonction localement constante qui donne la
caractéristique d'Euler--Poincaré des faisceaux inversibles paramétrés, et où
$N$ est un faisceau inversible quelconque sur $X \times_S P$ représentant le
point universel, $N_g = g_P^{*}N$, $N_{g'} = g'^{*}_P N$ ; le produit
$N_g N_{g'}^{-1}$ ne dépend pas du choix de $N$.\footnote{La notation $(*)$ de
la page est renommée $(*_{b})$ pour ne pas la confondre avec l'isomorphisme
$(*)$ de la page 11. Une note encadrée de la marge, rattachée par une flèche au
paragraphe suivant, dit l'essentiel : $N_g N_{g'}^{-1}$ est algébriquement
équivalent à zéro sur $P$ relativement à $S$, donc $M_g$ et $M_{g'}$ sont
algébriquement équivalents.}

Comme $P^{\delta}$ est un $B$-torseur, ses fibres sont connexes et $\chi$ y est
constante : on note $\chi(\delta)$ sa valeur. Sur $P^{\delta}$, le faisceau
$N_g N_{g'}^{-1}$ appartient à $\mathrm{Pic}^{0}$ ; c'est le point
$\psi(g') - \psi(g)$ de $A$, au sens du torseur d'Albanese des pages 22 et 23
(voir plus bas). Ainsi $M_g$ et $M_{g'}$ ont la même classe dans
$\mathrm{NS}_{P^{\delta}/S}(S)$.

\pagerange{2}{5}
\section*{L'application $\varphi$ et la section canonique $L(\delta)$ (pages 2 à 5)}

\subsection*{L'application $\varphi$}

Les translations de $B$ agissent trivialement sur $\mathrm{Pic}^{0}_{B/S}$ et
sur $\mathrm{NS}_{B/S}$ ; d'où des isomorphismes canoniques
$\mathrm{Pic}^{0}_{P^{\delta}/S} \simeq \mathrm{Pic}^{0}_{B/S}$ et
$\mathrm{NS}_{P^{\delta}/S} \simeq \mathrm{NS}_{B/S}$, qui ne dépendent pas du
point de $P^{\delta}$ choisi pour identifier $P^{\delta}$ à $B$. La classe de
$M_g|_{P^{\delta}}$ définit donc une section
\[
\varphi(\delta) \in \mathrm{NS}_{B/S}(S),
\]
indépendante de $g$ d'après $(*_{b})$. Elle se définit donc par descente sans
supposer l'existence d'une section — $X/S$ en acquiert une après le changement
de base fidèlement plat $X \to S$ —, et comme sa formation commute au
changement de base, on obtient un morphisme de faisceaux
\[
\varphi : \mathrm{NS}_{X/S} \longrightarrow \mathrm{NS}_{B/S}.
\]
Ce n'est pas un homomorphisme, et la page le souligne ; le calcul des pages 6
et 7 montre que c'est une application polynomiale, de degré $n - 1$ si $n$ est
la dimension relative de $X$.\footnote{La marge porte « détailler » en face de
l'argument de descente.}

\subsection*{Tordre $\mathrm{Pic}_{B/S}$ par le torseur $P^{\delta}$}

Soit $\lambda$ une section de $\mathrm{NS}_{B/S}$, et $\tilde{\lambda} : B \to
A$ l'homomorphisme qu'elle définit. Si $N$ est un faisceau inversible de classe
$\lambda$ sur $P^{\delta}$, chaque point $p$ de $P^{\delta}$ donne un
isomorphisme $B \to P^{\delta}$, $x \mapsto p + x$, et donc un faisceau $N_p$
sur $B$ ; on a $N_{p+b} = t_b^{*}N_p = N_p \otimes \tilde{\lambda}(b)$. Ainsi
$p \mapsto N_p$ est un morphisme $P^{\delta} \to \mathrm{Pic}^{\lambda}_{B/S}$
équivariant le long de $\tilde{\lambda}$, c'est-à-dire un morphisme de
$A$-torseurs $\tilde{\lambda}_{*}P^{\delta} \to \mathrm{Pic}^{\lambda}_{B/S}$,
et l'on obtient
\[
\mathrm{Pic}^{\lambda}_{P^{\delta}/S} \simeq
\mathrm{Pic}^{\lambda}_{B/S} \times^{A} \bigl(\tilde{\lambda}_{*}P^{\delta}\bigr)^{(-1)} .
\]
\footnote{La page écrit le second facteur sans l'inverse :
$\mathrm{Pic}^{\lambda}_{B/S} \overset{A}{\times} (P^{\delta}
\overset{B}{\times} A)$. Avec la convention de Mumford pour $\tilde{\lambda}$,
c'est l'inverse qui convient ; avec la convention opposée, la forme de la page.
L'inverse est retenu partout dans ce qui suit, parce que c'est lui qui donne à
la page 11 l'isomorphisme $(*)$ exactement tel que la page l'encadre, et parce
qu'il s'accorde avec $(**)$ de la page 18. On ne le signale plus à chaque
occurrence.}

Appliquons cela à $\lambda = \varphi(\delta)$ et au fibré $M_g|_{P^{\delta}}$ :
chaque section $g$ de $X$ donne un point $\ell(g)$ du $A$-torseur
\[
T_{\delta} = \mathrm{Pic}^{\varphi(\delta)}_{B/S} \times^{A}
\bigl(\widetilde{\varphi(\delta)}_{*}P^{\delta}\bigr)^{(-1)} .
\]
\footnote{La page 3 écrit ici « $L_g \mid P^{\delta}$ » ; c'est
$M_g|_{P^{\delta}}$ qui a la classe $\varphi(\delta)$, $\mathcal{L}_g$ vivant
sur $X \times_S P$.} La formation de $\ell(g)$ commutant au changement de base,
on l'applique au point universel — la base $X$, la section diagonale — et l'on
obtient un morphisme de schémas
\[
\ell_{\delta} : X \longrightarrow T_{\delta}.
\]

\subsection*{La section canonique}

Un morphisme de $X$ dans un torseur sous un schéma abélien se factorise de
façon unique par $\psi : X \to \mathrm{Alb}^{1}_{X/S}$ (c'est la propriété
universelle énoncée page 22), à travers un morphisme de torseurs compatible
avec un homomorphisme $u_{\delta} : A \to A$ des groupes structuraux. Or
$(*_{b})$ dit exactement comment $\ell(g)$ varie avec $g$ :
$\ell(g') = \ell(g) - \chi(\delta)\,\bigl(\psi(g') - \psi(g)\bigr)$. Donc
\[
u_{\delta} = -\chi(\delta)\,\mathrm{id}_A ,
\]
et la factorisation est un isomorphisme de $A$-torseurs
$\mathrm{Alb}^{-\chi(\delta)}_{X/S} \xrightarrow{\ \sim\ } T_{\delta}$,
c'est-à-dire une section canonique
\[
\boxed{\;L(\delta) \in \Gamma\bigl(S,\ \mathrm{Alb}^{\chi(\delta)}_{X/S}
\times^{A} \mathrm{Pic}^{\varphi(\delta)}_{B/S} \times^{A}
\bigl(\widetilde{\varphi(\delta)}_{*}P^{\delta}\bigr)^{(-1)}\bigr).\;}
\]
En particulier ce $A$-torseur est trivial, avec une trivialisation
canonique.\footnote{La transcription relève que l'exposant de $\mathrm{Alb}$
est $-\chi(\delta)$ dans la flèche et $\chi(\delta)$ dans le cadre. Les deux
sont justes ensemble : un morphisme de torseurs de source
$\mathrm{Alb}^{-\chi}$ est une section de $T \times^{A} \mathrm{Alb}^{\chi}$.
Le signe de $u_{\delta}$, sur lequel la marge écrit « vérifier si le signe est
correct ! », est bien $-\chi(\delta)$ avec la convention de signe de la
page 22 ; la page 11 le retrouve par un autre chemin.}

\pagerange{5}{7}
\section*{Le calcul de $\varphi(\delta)$ par Riemann--Roch (pages 5 à 7)}

Pour calculer $\chi(\delta)$ et $\varphi(\delta)$ on se ramène à
$S = \operatorname{Spec} k$, $k$ algébriquement clos. Un point $p$ de
$P^{\delta}(k)$, c'est-à-dire un faisceau inversible $L$ de classe $\delta$,
identifie $P^{\delta}$ à $B$, et $\mathcal{L}_g$ restreint à $X \times
P^{\delta}$ devient $\mathcal{L}^{0} \otimes \mathrm{pr}_1^{*}L$, où
$\mathcal{L}^{0}$ est le faisceau de Poincaré sur $X \times B$, normalisé en
$g$. Il s'agit donc de calculer la classe de $\det R\mathrm{pr}_{2*}
(\mathcal{L}^{0} \otimes \mathrm{pr}_1^{*}L)$ modulo équivalence algébrique.
Comme $\mathrm{NS}(B)$ est sans torsion, il suffit de le faire dans
$\mathrm{NS}(B) \otimes \mathbb{Q}$, c'est-à-dire de calculer la partie de
degré $1$ du caractère de Chern.\footnote{Cette justification est celle d'une
phrase biffée de la page 6 ; la phrase gardée dit seulement « mod équivalence
algébrique \ldots\ dans $\mathrm{NS}(B/S) \otimes \mathbb{Q}$ ».}

Si $X$ est lisse de dimension $n$ — la page dit aussi « intersection
complète » —, Grothendieck--Riemann--Roch donne
\[
\varphi(\delta) = \Bigl[\mathrm{pr}_{2*}\bigl(e^{D}\cdot
\mathrm{pr}_1^{*}(\mathrm{Td}_X\, e^{\delta})\bigr)\Bigr]_{1},
\qquad D = c_1(\mathcal{L}^{0}),
\]
où $[\ ]_1$ désigne la composante de degré $1$.\footnote{La page note
« $\Delta = \mathrm{cl}\,\Delta$ », là où l'on attend la classe de $L$ ; on
écrit $\delta$. Elle écrit aussi le résultat sous la forme $[\exp D](T_{X/k}
\exp \Delta)$, avec un « $\exp D$ » ajouté au crayon : c'est l'opération par
correspondance de la classe $e^{D}$.} On travaille en cohomologie
$\ell$-adique, $\ell$ inversible dans $k$, en choisissant un isomorphisme
$\mathbb{Z}_\ell(1) \simeq \mathbb{Z}_\ell$. La normalisation de
$\mathcal{L}^{0}$ le long de $g \times B$ et de $X \times 0$ fait que $D$
appartient à la composante de Künneth $H^{1}(X) \otimes H^{1}(B)$, où il
réalise la dualité entre $H^{1}(B)$ et $H^{1}(X)$. Écrivons
$\mathrm{Td}_X\, e^{\delta} = \sum_{i} a_i$, avec $a_i \in H^{2i}(X)$. Le terme
$\mathrm{pr}_1^{*}(a_i)\,D^{j}/j!$ a le degré partiel $2i + j$ en $X$ ; son
image par $\mathrm{pr}_{2*}$ est nulle sauf si $2i + j = 2n$, et elle est alors
de degré $j$ sur $B$. La composante de degré $2$ (la classe $c_1$) demande
$j = 2$, $i = n - 1$ :
\[
\boxed{\;\varphi(\delta) = \tfrac{1}{2}\,\mathrm{pr}_{2*}\bigl(\mathrm{pr}_1^{*}(a_{n-1})\, D^{2}\bigr),
\qquad a_{n-1} = \Bigl[\mathrm{Td}_X\, e^{\delta}\Bigr]_{n-1}.\;}
\]
Comme $a_{n-1} = \sum_{k} \mathrm{Td}_{n-1-k}(X)\,\delta^{k}/k!$,
$\varphi$ est un polynôme de degré $n - 1$ en $\delta$.

\pagerange{7}{9}
\section*{Premier exemple : les courbes (pages 7 à 9)}

Pour $n = 1$, $\delta$ est le degré $d$, $a_0 = 1$ et
$\varphi(\delta) = \tfrac{1}{2}\mathrm{pr}_{2*}(D^{2})$, indépendant de $d$. Dans
une base symplectique $(e_1, f_1, \ldots, e_g, f_g)$ de $H^{1}(X)$, la page
calcule $\tfrac{1}{2}D^{2} = \eta \otimes \sum e_i f_i$, d'où
$\varphi(\delta) = \pm \sum e_i \wedge f_i$ : au signe près, la classe de la
polarisation canonique $\theta$ de la jacobienne $B$. Le signe est
\[
\varphi(\delta) = -\theta .
\]
\footnote{La page écrit $+\sum e_i \wedge f_i$, et conclut : « c'est la classe
de polarisation canonique de la jacobienne ». Le signe dépend de
l'identification de $H^{1}(X)$ au dual de $H^{1}(B)$, que la page ne fixe pas.
Il est intrinsèque une fois $\varphi(\delta)$ défini comme classe d'un fibré
sur $B$, et c'est $-\theta$ : pour $d$ grand, $R\mathrm{pr}_{2*}$ est le
\emph{fibré de Picard}, dont la classe de Chern totale est $e^{-\theta}$
(Mattuck, 1961) ; et sur une courbe elliptique, avec $d = 1$, on trouve
directement $\mathrm{pr}_{2*}\mathcal{L}_g \simeq \mathcal{O}(-g)$, de degré
$-1$. Ce signe est d'ailleurs celui qu'exige la page 10 elle-même, où c'est
$-\varphi(\delta)$ qui est une polarisation.} Riemann--Roch donne
$\chi(\delta) = 1 - g + d$.\footnote{La marge porte ici « vérifier que le
signe est correct ». Une ligne biffée de la page 7 écrit
$\mathrm{Td}_X = 1 - K/2$, ce qui est juste.}

La page traduit ensuite les trois facteurs du torseur de $L(\delta)$ en
$B$-torseurs, par l'isomorphisme $\phi_{\Theta} : B \to A$ de la polarisation
thêta. Sous $\phi_{\Theta}$, $\mathrm{Alb}^{1}_{X/S}$ devient
$P^{-1} = \mathrm{Pic}^{-1}_{X/S}$ (« attention au signe ! » — c'est ce que
confirme la page 23), donc $\mathrm{Alb}^{\chi(\delta)}$ devient
$P^{g-1-d}$ ; le facteur poussé par $\widetilde{\varphi(\delta)}$ redevient
$P^{\delta} = P^{d}$ ; et $\mathrm{Pic}^{\varphi(\delta)}_{B/S}$ s'identifie,
par le diviseur $\Theta$, à une puissance de $P^{g-1}$. La page trouve pour le
produit
\[
P^{2g-2} = \mathrm{Pic}^{2g-2}_{X/S},
\]
et observe que ce n'est « pas très excitant » : on s'attend à ce que la section
canonique soit la classe $K$ du fibré canonique, ce qui reste « à
prouver ! ».\footnote{Ce décompte est le point le moins sûr du dossier. Les
exposants sont $\varepsilon_1(g - 1 - d) + \varepsilon_2(g-1) +
\varepsilon_3 d$, avec trois signes qui dépendent de conventions que la page
signale elle-même comme douteuses. Que les termes en $d$ se compensent est
nécessaire — la section existe pour tout $d$, et sur la courbe universelle
$\mathrm{Pic}^{m}$ n'a de section que si $2g - 2$ divise $m$ (voir la note
suivante) — et ils se compensent bien. Mais avec les conventions fixées ici
($\varphi(\delta) = -\theta$, le torseur inverse de la page 3, et
$\mathrm{Pic}^{\theta}_{B/S} \simeq P^{g-1}$ par $\xi \mapsto
t_{\xi}^{*}W_{g-1}$), notre décompte donne $\varepsilon_2 = -1$ et un
exposant total nul, non $2g - 2$. Nous ne l'avons pas vérifié autrement ; le
résultat de la page n'est donc pas confirmé, et le nôtre pas davantage. La
section de $\mathrm{Pic}^{2g-2}$ égale à $K$ que la page attend ici est, elle,
bien obtenue page 27, par un autre chemin.}

\emph{N.B.} (page 9). Pour bien faire, dit la page, il faudrait montrer que
sur l'espace de modules $S$ des courbes de genre $g \geq 1$ (avec un point
marqué si $g = 1$), le groupe des sections de $\mathrm{Pic}_{X/S}$ de la
courbe universelle est engendré par $K$ — en particulier que $\mathrm{Pic}^{0}$
n'a que la section nulle — et qu'il suffirait de le faire en caractéristique
nulle.\footnote{C'est la \emph{conjecture de Franchetta} (1954), sous sa forme
forte. Elle est démontrée : en caractéristique nulle par Arbarello et Cornalba
(1987) et d'autres, en toute caractéristique par Schröer (2003). La page ne
cite personne ; une insertion d'une ligne au-dessus de « courbes » est
illisible.}

\pagerange{9}{13}
\section*{Second exemple : les schémas abéliens (pages 9 à 13)}

\subsection*{La transformation de Fourier et l'inversion des polarisations}

Soit maintenant $X = A$ un schéma abélien de dimension relative $n$, de sorte
que $B$ est son dual et que $\mathrm{Pic}^{0}_{B/S}$ redevient $A$. Ici
$\mathrm{Td}_A = 1$, donc $a_{n-1} = \delta^{n-1}/(n-1)!$. La correspondance
$e^{D}$ est la \emph{transformation de Fourier}
\[
\Phi : \mathrm{CH}(A)_{\mathbb{Q}} \xrightarrow{\ \sim\ } \mathrm{CH}(B)_{\mathbb{Q}},
\]
qui échange produit d'intersection et produit de Pontryagin et, en
cohomologie, envoie le degré $2i$ sur le degré $2(n - i)$. On obtient
\[
\varphi(\delta) = \frac{1}{(n-1)!}\,\Phi\bigl(\delta^{n-1}\bigr)
= \pm \frac{1}{(n-1)!}\,\Phi(\delta)^{*(n-1)},
\]
où $\Phi(\delta)$ est une classe de courbe. Ainsi $\varphi$ est une
application polynomiale de degré $n - 1$
de $\mathrm{NS}_{A/S}$ dans $\mathrm{NS}_{B/S}$.\footnote{La page écrit la
seconde forme, avec le signe $+$. Avec la normalisation de Beauville,
$\Phi(xy) = (-1)^{n}\Phi(x) * \Phi(y)$, d'où le signe $(-1)^{n}$ entre les
deux expressions ; la première, qui n'en a pas, est celle qui résulte
directement de la page 7. La transformation de Fourier des cycles porte
aujourd'hui les noms de Mukai (1981, pour les catégories dérivées) et de
Beauville (1983, pour les groupes de Chow) ; sa version cohomologique est plus
ancienne. La page dit « on sait », sans référence.}

Le fait central est la relation encadrée de la page 9 :
\[
\boxed{\;\widetilde{\varphi(\delta)} \circ \tilde{\delta} = -\chi(\delta)\,\mathrm{id}_A,
\qquad
\tilde{\delta} \circ \widetilde{\varphi(\delta)} = -\chi(\delta)\,\mathrm{id}_B .\;}
\]
La page la démontre au paragraphe « Remarque » de la page 11 : la section
$L(\delta)$ est canonique, donc invariante par les automorphismes de
$(X/S, \delta)$, en particulier par les translations de $X = A$. Une
translation $t_a$ agit sur le torseur de $L(\delta)$ par
$\chi(\delta)\,a + \widetilde{\varphi(\delta)}\,\tilde{\delta}(a)$ — le premier
terme vient de $\mathrm{Alb}^{\chi(\delta)}$, le second du facteur
$P^{\delta}$, le facteur $\mathrm{Pic}_{B/S}$ étant fixe —, et cette
translation doit être nulle pour tout $a$.\footnote{La marge : « Attention aux
signes ?? », et page 9, « Faire calcul direct, en faisant attention aux
signes ! ». Le signe est juste. En langage d'aujourd'hui, c'est la
proposition de Mukai (1981) : pour $L$ ample, $\varphi_L^{*}\widehat{L} \simeq
H^{0}(L) \otimes L^{-1}$, dont le déterminant donne
$\varphi_L^{*}\varphi(\delta) = -\chi(\delta)\,\delta$, d'où la première
égalité pour $\delta$ ample, puis pour tout $\delta$ par densité du cône
ample, les deux membres étant polynomiaux en $\delta$.}

Conséquences (page 10). Si $\chi(\delta) \neq 0$, $\varphi(\delta)$ est non
dégénérée ; si $\delta$ est une polarisation, $\delta' := -\varphi(\delta)$ en
est une sur $B$. Pour $\delta$ ample, $\deg \tilde{\delta} = \chi(\delta)^{2}$,
$\deg \tilde{\delta}' = \chi(\delta')^{2}$ et $\deg
\tilde{\delta}'\tilde{\delta} = \chi(\delta)^{2n}$ donnent
\[
\chi(\delta') = \chi(\delta)^{n-1},
\]
et l'égalité vaut pour tout $\delta$, les deux membres étant des polynômes en
$\delta$.\footnote{L'argument des degrés ne donne $\chi(\delta')$ qu'au signe
près ; c'est la positivité dans le cas ample, puis la densité du cône ample,
qui le fixe. La réciproque que la page énonce ensuite — $\varphi(\delta)$ non
dégénérée entraîne $\chi(\delta) \neq 0$ — suppose $n \geq 2$ : pour une
courbe elliptique, $\varphi(\delta)$ est toujours non dégénérée, même en degré
nul. La marge oblique, en partie illisible, porte « $\delta' =
-\varphi(\delta)$ ; je pense ».} L'inversion $\delta \mapsto \delta'$ est une
involution au scalaire près :
\[
\delta'' = -\varphi(\delta') = \chi(\delta)^{n-2}\,\delta,
\qquad\text{donc}\qquad
\varphi(\varphi(\delta)) = (-1)^{n}\,\chi(\delta)^{n-2}\,\delta .
\]
\footnote{La page écrit « si $\delta' = \varphi(\delta)$ », puis conclut
$\varphi(\varphi(\delta)) = \delta\,\chi(\delta)^{n-2}$ sans signe, ce qui
n'est juste que pour $n$ pair ; la note de la transcription dit que les
signes y sont ceux de la page. La relation sur $\delta''$, qui est celle dont
la page 12 se sert, est juste telle quelle.} Pour $n = 1$, une courbe
elliptique, on trouve que $\varphi(\delta)$ est la même pour tout $\delta$ :
c'est, d'après les pages 7 et 8, l'opposée $-\theta$ de la polarisation
canonique.\footnote{La page écrit « on trouve $\delta\,\chi(\delta)$ », là où
la formule donne $\delta\,\chi(\delta)^{-1}$, puis « mieux, pour tout
$\delta$, $\varphi(\delta)$ est la polarisation canonique » ; le signe est
celui de la note de la page 8. Pour $n \geq 2$, conclut-elle, la formule
« inspire confiance ».}

\emph{N.B.} (page 10). La page s'attend à ce que, sur le site modulaire des
schémas abéliens polarisés de degré $\chi(\delta)$ donné, $\delta'$ ne soit pas
divisible, et même engendre les sections de $\mathrm{NS}_{B/S}$ : $\delta'$
serait l'« inversion » la meilleure possible d'une polarisation.\footnote{Telle
quelle, l'attente est fausse dès qu'un facteur élémentaire autre que le dernier
dépasse $1$. Pour une polarisation de type $(d_1, \ldots, d_n)$,
$\tilde{\delta}'\tilde{\delta} = \chi(\delta) = d_1 \cdots d_n$, tandis que la
\emph{polarisation duale} $\delta^{\vee}$ (Birkenhake--Lange, § 14.4) vérifie
$\tilde{\delta}^{\vee}\tilde{\delta} = d_n$ ; donc $\delta' = d_1 \cdots
d_{n-1}\,\delta^{\vee}$, divisible par exemple par $2$ pour le type $(2,2)$.
C'est $\delta^{\vee}$, non $\delta'$, qui est « ce qu'on peut faire de mieux »,
et les deux coïncident exactement pour les types $(1, \ldots, 1, d)$. Le reste
de la note est en partie illisible.}

\subsection*{L'isomorphisme canonique $(*)$ (page 11)}

Pour $X = A$, le torseur $\mathrm{Alb}^{1}_{A/S}$ est trivialisé par la section
unité, et la section $L(\delta)$ devient un isomorphisme canonique de
$A$-torseurs
\[
(*) \qquad
\boxed{\;\mathrm{Pic}^{\delta'}_{B/S} \simeq \tilde{\delta}'_{*}\,\mathrm{Pic}^{\delta}_{A/S}\;}
\]
— le $B$-torseur $\mathrm{Pic}^{\delta}_{A/S}$ poussé en un $A$-torseur par
$\tilde{\delta}' : B \to A$. Il a une description directe : c'est
$L \mapsto (\det \widehat{L})^{-1}$, où $\widehat{L} = R\mathrm{pr}_{2*}
(\mathcal{P} \otimes \mathrm{pr}_1^{*}L)$ est la transformée de Fourier--Mukai
de $L$ ; l'équivariance le long de $\tilde{\delta}'$ résulte de
$\widehat{L \otimes \mathcal{P}_{\beta}} = t_{\beta}^{*}\widehat{L}$ et du
théorème du carré.\footnote{Dans le cadre, l'exposant $\delta$ de
$\mathrm{Pic}_{A/S}$ est écrit par-dessus un premier exposant, peut-être
$-\delta$, et la longue note oblique de la marge, en grande partie illisible,
remarque que $\mathrm{Pic}^{-\delta}_{A/S}$ et $\mathrm{Pic}^{\delta}_{A/S}$
s'identifient canoniquement et demande « à vérifier ». Avec la convention de
la page 3 fixée ici, c'est l'exposant $\delta$ qui est juste. La description
par Fourier--Mukai est nôtre ; c'est elle que la page 19 retrouve en
substance.}

\emph{Exemple} (page 11). Si $\delta$ est principale, $\tilde{\delta}'$ est
l'inverse de $\tilde{\delta}$ et $\delta'$ est la polarisation transportée. La
page lit alors $(*)$ comme un isomorphisme $\mathrm{Pic}^{\delta}_{A/S} \simeq
\mathrm{Pic}^{-\delta}_{A/S}$ et commence de le comparer à l'isomorphisme
défini par la symétrie $-1$ ; le passage est biffé.\footnote{Un bloc de quatre
lignes encadré et biffé de grandes boucles, dont on ne lit que des bribes. Avec
les conventions fixées ici, transporté par $\tilde{\delta}$, $(*)$ est
l'identité de $\mathrm{Pic}^{\delta}_{A/S}$, puisque $\varphi_L^{*}\widehat{L}
\simeq H^{0}(L) \otimes L^{-1}$ avec $H^{0}(L)$ de rang $1$.}

\subsection*{Puissances de torseurs canoniquement triviales (pages 12 et 13)}

\emph{Remarque} (page 12). Remplaçons $\delta$ par $\lambda\delta$, $\lambda \in \mathbb{Z}$. Comme
$\varphi$ est homogène de degré $n - 1$, $(\lambda\delta)' =
\lambda^{n-1}\delta'$, et $(*)$ pour $\lambda\delta$ s'écrit
\[
\bigl(\mathrm{Pic}^{\delta'}_{B/S}\bigr)^{(\lambda^{n-1})} \simeq
\bigl(\tilde{\delta}'_{*}\mathrm{Pic}^{\delta}_{A/S}\bigr)^{(\lambda^{n})}
\simeq \bigl(\mathrm{Pic}^{\delta'}_{B/S}\bigr)^{(\lambda^{n})},
\]
d'où une trivialisation de $(\mathrm{Pic}^{\delta'}_{B/S})^{(\lambda^{n} -
\lambda^{n-1})}$ ; pour $\lambda = -1$ et $n \geq 2$, une trivialisation de
$(\mathrm{Pic}^{\delta'}_{B/S})^{(2)}$, et $2$ est le p.g.c.d. des
$\lambda^{n} - \lambda^{n-1}$.\footnote{Les exposants en $\lambda$ sont tracés
très vite et lus sans certitude ; la page écrit $\mathrm{Pic}^{-\lambda\delta}$
au second membre, ce qui suit la forme biffée de $(*)$. Ces trivialisations
sont canoniques à condition que les isomorphismes $(*)$ pour $\delta$ et pour
$\lambda\delta$ soient compatibles aux puissances : c'est ce que la marge
demande de vérifier, et le dossier ne le fait pas. Pour $n = 1$ le p.g.c.d. est
$1$.} En échangeant les rôles de $A$ et $B$ grâce à $\delta'' =
\chi(\delta)^{n-2}\delta$, on obtient, pour $n \geq 2$, une trivialisation de
$(\mathrm{Pic}^{\delta}_{A/S})^{(2\chi(\delta)^{n-2})}$.

La page 13 refait le compte plus finement. En composant $(*)$ pour $\delta$ et
pour $\delta'$, et comme $\tilde{\delta}''\tilde{\delta}' = \chi(\delta')
= \chi(\delta)^{n-1}$,
\[
\bigl(\mathrm{Pic}^{\delta}_{A/S}\bigr)^{(\chi(\delta)^{n-2})} \simeq
\mathrm{Pic}^{\delta''}_{A/S} \simeq
\bigl(\mathrm{Pic}^{\delta}_{A/S}\bigr)^{(\chi(\delta)^{n-1})},
\]
d'où la trivialité de $(\mathrm{Pic}^{\delta}_{A/S})^{(\chi^{n-2}(\chi - 1))}$,
et, en remplaçant $\delta$ par $\lambda\delta$, de la puissance
$\chi^{n-2}(\lambda^{n}\chi - \lambda^{n-1})$. Avec $\lambda = \pm 1$ on
retrouve l'exposant $2\chi(\delta)^{n-2}$ ; et si $\chi(\delta)$ est pair,
$\chi - 1$ est impair, ce qui donne la trivialité de
$(\mathrm{Pic}^{\delta}_{A/S})^{(\chi(\delta)^{n-2})}$. \emph{Exemple} (page 13), $n = 2$ :
$(\mathrm{Pic}^{\delta}_{A/S})^{(2)}$ est trivial, et
$\mathrm{Pic}^{\delta}_{A/S}$ lui-même l'est si $\chi(\delta)$ est
pair.\footnote{La page indice ici $\mathrm{Pic}^{\delta}$ par $B/S$, là où les
lignes précédentes ont $A/S$ ; c'est $A/S$. Une ligne cernée et en partie
barrée parle d'un « p.p.c.m. » (lecture douteuse) : c'est du p.g.c.d. des
exposants qu'il s'agit, comme le dit la ligne biffée qui précède. Deux remarques
pour situer ces énoncés. D'abord, la trivialité de
$(\mathrm{Pic}^{\delta}_{A/S})^{(2)}$ est élémentaire pour tout $\delta$ :
$L \mapsto L \otimes [-1]^{*}L$ est une section de
$\mathrm{Pic}^{2\delta}_{A/S}$ qui ne dépend pas de $L$. Ce que le dossier
apporterait, c'est l'énoncé de parité, et l'identification de ses
trivialisations aux trivialisations canoniques que les marges demandent de
vérifier. Ensuite, l'énoncé pour $n = 2$ et $\chi(\delta)$ pair dépend de ces
compatibilités, que le dossier n'établit pas ; cette lecture n'a pas pu le
confirmer.} Une ligne biffée en tête de la page 13 annonçait « g)
Généralisation des constructions précédentes » ; c'est le paragraphe qui
suit.

\pagerange{14}{16}
\section*{Généralisation : un complexe parfait (pages 14 à 16)}

On revient au cadre général, et l'on se donne un complexe parfait
$\mathcal{F}$ sur $X$ (ou seulement parfait relativement à $S$). On pose
$M_g^{\mathcal{F}} = \det R\mathrm{pr}_{2*}(\mathrm{pr}_1^{*}\mathcal{F} \otimes
\mathcal{L}_g)$. Tout ce qui précède se refait mot pour mot, $\chi$ étant
remplacé par $\chi_{\mathcal{F}}$, le rang de $R\mathrm{pr}_{2*}
(\mathrm{pr}_1^{*}\mathcal{F} \otimes \mathcal{L}_g)$ : on obtient
$\varphi_{\mathcal{F}} : \mathrm{NS}_{X/S} \to \mathrm{NS}_{B/S}$, additive en
$\mathcal{F}$ — elle ne dépend que de la classe de $\mathcal{F}$ dans le groupe
de Grothendieck de $D_{\mathrm{parf}}(X)$, le déterminant étant multiplicatif
sur les triangles —, un morphisme $\ell^{\mathcal{F}}_{\delta}$ et une section
canonique
\[
L_{\mathcal{F}}(\delta) \in \Gamma\Bigl(S,\ \mathrm{Alb}^{\chi_{\mathcal{F}}(\delta)}_{X/S}
\times^{A} \mathrm{Pic}^{\varphi_{\mathcal{F}}(\delta)}_{B/S} \times^{A}
\bigl(\widetilde{\varphi_{\mathcal{F}}(\delta)}_{*}P^{\delta}\bigr)^{(-1)}\Bigr).
\]
\footnote{La page écrit le troisième facteur avec
$\mathrm{Pic}^{\delta}_{X/S}$, qui est $P^{\delta}$, et un des produits sans
l'indice $A$ ; l'exposant de $\mathrm{Alb}$ change de signe d'une ligne à
l'autre, comme page 4 et pour la même raison.} Le calcul se ramène encore au
cas d'un corps, et $\delta$ y disparaît, à condition de remplacer
$\mathcal{F}$ par $L_{\delta} \otimes \mathcal{F}$. Par Riemann--Roch,
\[
\chi_{\mathcal{F}}(\delta) = \int_X \mathrm{ch}(L_{\delta} \otimes \mathcal{F})\,\mathrm{Td}_X,
\qquad
\varphi_{\mathcal{F}}(\delta) = \Bigl[\mathrm{pr}_{2*}\bigl(e^{D}\cdot
\mathrm{pr}_1^{*}(\mathrm{ch}(L_{\delta} \otimes \mathcal{F})\,\mathrm{Td}_X)\bigr)\Bigr]_{1}.
\]

\emph{Ex. 1} (page 15), $X$ une courbe. Si $\mathcal{F}$ est de rang $\rho$ et
de déterminant de degré $c$, on a $\mathrm{ch}(L_{\delta} \otimes \mathcal{F})
= \rho + (\rho\delta + \det\mathcal{F})$ et $\mathrm{Td}_X = 1 - K/2$, d'où
\[
\chi_{\mathcal{F}}(\delta) = \rho\,(1 - g + d) + c,
\qquad
\varphi_{\mathcal{F}}(\delta) = \rho\,\varphi(\delta) = -\rho\,\theta .
\]
\footnote{La page fait le calcul avec $\rho = 1$, et trouve
$\chi_{\mathcal{F}}(\delta) = 1 - g + d + c$ et $\varphi_{\mathcal{F}}(\delta)
= \varphi(\delta)$ ; elle écrit aussi $\mathrm{Todd}(X) = 1 + \tfrac{1}{2}K$,
tout en notant sous ce terme « degré $1 - g$ », qui est le degré de $-K/2$. Le
signe de $\varphi(\delta)$ est celui de la page 8. Les lettres $\rho$ et
$\gamma$ (pour $\det \mathcal{F}$) sont lues sans certitude.} En comparant
$L_{\mathcal{F}}(\delta)$ à la puissance $\rho$-ième de $L(\delta)$, on obtient
une section canonique de $\mathrm{Alb}^{c}_{X/S} \simeq \mathrm{Pic}^{c}_{X/S}$,
et la page conclut : « Il n'y a pas de miracle ici », car c'est simplement la
classe de $\det\mathcal{F}$.

\emph{Ex. 2} (page 16), $X = A$ un schéma abélien. Avec $\mathrm{ch}(L_{\delta}
\otimes \mathcal{F}) = \sum_{i=0}^{n} a_i$, on a $\chi_{\mathcal{F}}(\delta) =
\deg a_n$ et $\varphi_{\mathcal{F}}(\delta) = \tfrac{1}{2}\mathrm{pr}_{2*}
(\mathrm{pr}_1^{*}(a_{n-1})\,D^{2})$, la composante de degré $1$ de
$\Phi(a_{n-1})$.\footnote{La page fait commencer la somme à $i = 1$.}

\pagerange{16}{20}
\section*{Seconde interprétation : le déterminant d'une transformation de Fourier (pages 16 à 20)}

Le second paragraphe f) reconstruit $\ell_{\delta}$ « en intégrant
exclusivement sur $X \times_S B$ ». Soit $g$ une section de $X$ et $\lambda$
une section de $\mathrm{Pic}_{X/S}$, représentée grâce à $g$ par un faisceau
inversible $L = L(\lambda, g)$ trivialisé le long de $g$ ; soit $W(g)$ le
faisceau de Poincaré sur $X \times_S B$ normalisé le long de $g$. On pose
\[
\ell_g(\lambda) = \bigl[\det R\mathrm{pr}_{2*}\bigl(\mathrm{pr}_1^{*}L \otimes W(g)\bigr)\bigr]
\in \mathrm{Pic}_{B/S}(S),
\]
ce qui définit, par changement de base, un morphisme
\[
\ell' : X \times_S \mathrm{Pic}_{X/S} \longrightarrow \mathrm{Pic}_{B/S}.
\]
\footnote{La page écrit $\mathrm{pr}_{2*}$ sans $R$, et dit $\ell'$
« homomorphisme (additif en $\mathrm{Pic}_{X/S}$) ». Il n'est pas additif :
$\ell_g$ envoie $\mathrm{Pic}^{\delta}_{X/S}$ dans
$\mathrm{Pic}^{\varphi(\delta)}_{B/S}$ et $\varphi$ n'est pas additive. Ce
qui est vrai, c'est la formule $(**)$ ci-dessous, qui fait de $\ell_g$ une
application affine sur chaque $\mathrm{Pic}^{\delta}_{X/S}$. Dans le cadre,
l'indice de $\mathrm{Pic}$ se lit $X/S$ ou $X/B$.}

Le changement de section se calcule comme à la page 1 :
$L(\lambda, g') = L(\lambda, g) \otimes g'^{*}L(\lambda, g)^{-1}$ et
$W(g') = W(g) \otimes \mathrm{pr}_2^{*}\, g'^{*}_B W(g)^{-1}$, d'où
\[
\ell_{g'}(\lambda) = \ell_g(\lambda) - \chi(\lambda)\,\bigl(\psi_X(g') - \psi_X(g)\bigr),
\]
où $\psi_X : X \to \mathrm{Alb}^{1}_{X/S}$ est le morphisme canonique —
la même dépendance en $-\chi$ qu'à la page 4. On en déduit, pour chaque
section $\delta$ de $\mathrm{NS}_{X/S}$, un morphisme
$\ell_{\delta} : X \times_S \mathrm{Pic}^{\delta}_{X/S} \to
\mathrm{Pic}^{\varphi(\delta)}_{B/S}$.

La page 18 énonce alors, pour $\beta \in B(S)$ et $\lambda$ de classe $\delta$,
\[
(**) \qquad \ell_g(\lambda + \beta) = \ell_g(\lambda) + \widetilde{\varphi(\delta)}(\beta).
\]
\footnote{Le signe entre les deux termes est noyé sous une tache d'encre ; la
transcription lit $+$. Un paragraphe encadré et biffé de quatre diagonales
annonçait un autre argument, où il s'agissait de prouver que l'homomorphisme
associé était $-\widetilde{\varphi(\delta)}$.} La preuve est une propriété de
la transformation de Fourier--Mukai $\mathfrak{F}_{\mathcal{L}}$ : tordre par
$L_{\beta} \in \mathrm{Pic}^{0}$ revient à translater la transformée,
$\mathfrak{F}_{\mathcal{L}}(L \otimes L_{\beta}) =
t_{\beta}^{*}\,\mathfrak{F}_{\mathcal{L}}(L)$ — la page écrit le second membre
comme convolution avec la masse de Dirac $\varepsilon_{-\beta}$, ce qui est la
même chose —, et le déterminant d'un translaté se calcule par le théorème du
carré : $\det t_{\beta}^{*}E = \det E \otimes \widetilde{[\det E]}(\beta)$. Avec
$[\det E] = \varphi(\delta)$, c'est $(**)$.\footnote{La lettre gothique
$\mathfrak{F}_{\mathcal{L}}$ est lue sans certitude ; la page renvoie à des
« notations développées ailleurs ». La page 19 conclut avec
$\widetilde{\varphi(\delta)}(-\beta)$, ce qui contredit $(**)$ à moins de
changer de convention pour $\tilde{\lambda}$ ; avec celle de Mumford, fixée
ici, c'est $+\beta$, comme $(**)$. La marge : « vérif., avec \ldots\ signes
\ldots ».}

Par $(**)$, $\ell_g$ est un isomorphisme de $A$-torseurs
$\widetilde{\varphi(\delta)}_{*}\mathrm{Pic}^{\delta}_{X/S} \to
\mathrm{Pic}^{\varphi(\delta)}_{B/S}$, c'est-à-dire un point du torseur
$T_{\delta}$ de la page 3 ; la variation en $g$ en fait un morphisme
$X \to T_{\delta}$, qui est le $\ell_{\delta}$ du paragraphe d), « comme on
vérifie facilement ». C'est cette lecture qui justifie, après coup, la
description de $(*)$ par Fourier--Mukai donnée plus haut.

\emph{Remarque} (pages 19 et 20). Cette interprétation ne fait intervenir que
des faisceaux sur $B$, non sur $\mathrm{Pic}^{\delta}_{X/S}$. Elle admet une
variante sans section : on se donne un homomorphisme de groupes
$i : B \to \mathrm{Pic}_{X/S}$, ce qui définit un torseur $Q$ sous $A$ et un
morphisme canonique $\tau : X \to Q$ ; la construction donne alors, pour un
$B$-torseur $P$ muni d'un morphisme compatible
$P \to \mathrm{Pic}^{\delta}_{X/S}$, un morphisme $X \to
\mathrm{Pic}^{\varphi(\delta)}_{B/S} \times^{A}
(\widetilde{\varphi(\delta)}_{*}P)^{(-1)}$ qui se factorise par
$Q^{(-\chi(\delta))}$, donc une section canonique de
$Q^{(\chi(\delta))} \times^{A} \mathrm{Pic}^{\varphi(\delta)}_{B/S}
\times^{A} (\widetilde{\varphi(\delta)}_{*}P)^{(-1)}$. Une parenthèse y voit le
commencement d'une étude de la « petite topologie de Albanese ».\footnote{Le
début de la remarque est très incertain : on lit « si
$\mathrm{Pic}^{0}_{X/S}$ admet pas une section », et l'objet donné, au bas de
la page 19, est coupé par un mot illisible. L'étiquette $\tau$ de la flèche
est lue sans certitude, comme toute la parenthèse sur la « petite topologie
de Albanese ». Dans les trois dernières formules, $P$ est écrit par-dessus un
$\mathrm{Pic}^{\delta}_{X/S}$ couvert de hachures. Un paragraphe b) de la
remarque, au bas de la page 20, est biffé de boucles croisées et presque
illisible ; la page 21 est blanche.}

\pagerange{22}{24}
\section*{La reprise : le torseur d'Albanese et sa convention de signe (pages 22 à 24)}

La seconde suite reprend les hypothèses du début : $X/S$ propre, plat, de
présentation finie, $f_{*}\mathcal{O}_X = \mathcal{O}_S$ universellement,
$B = \mathrm{Pic}^{0}_{X/S}$ schéma abélien, $A = \mathrm{Pic}^{0}_{B/S}$. Elle
construit le $A$-torseur $\mathrm{Alb}^{1}_{X/S}$ et le morphisme
$\psi : X \to \mathrm{Alb}^{1}_{X/S}$ dont la première suite s'est servie.
Il revient au même de se donner le morphisme
\[
\Delta_{\psi} : X \times_S X \longrightarrow A,
\qquad \Delta_{\psi}(x, y) = \psi(y) - \psi(x),
\qquad \Delta_{\psi}(x, x) = 0,
\]
défini sur les sections par $\Delta_{\psi}(x, y) = N_y N_x^{-1}$ dans
$\mathrm{Pic}^{0}_{B/S}(S)$, où $N$ est un faisceau de Poincaré sur
$X \times_S B$.\footnote{La page note $\delta$ cette application et l'appelle
« homomorphisme » ; c'est un morphisme de schémas. Le faisceau $N$ existe dès
que $X/S$ a une section, et le cas général s'obtient par descente.} La page
annonce la propriété universelle de $(\mathrm{Alb}^{1}_{X/S}, \psi)$ pour les
morphismes de $X$ dans les torseurs sous des schémas abéliens, et demande si
elle vaut pour des schémas en groupes commutatifs quelconques.\footnote{La page
dit « on prouve », sans preuve. Sur un corps algébriquement clos la réponse à
la question est oui : le sous-groupe engendré par l'image d'une variété
complète et connexe est complet, donc une variété abélienne. Au-dessus d'une
base quelconque, le dossier laisse la question ouverte, et nous aussi.}

Le couple $(\mathrm{Alb}^{1}, \psi)$ n'est déterminé qu'au signe près :
$(Q^{(-1)}, -\psi)$ est une autre solution du même problème universel,
canoniquement isomorphe à la première, de sorte que l'identification du groupe
des translations de $\mathrm{Alb}^{1}$ à $\mathrm{Pic}^{0}_{B/S}$ demande une
\emph{convention de signe}. La page adopte la première, celle de la formule
ci-dessus ; c'est celle qu'on a suivie partout.

Pour $X$ lisse de dimension relative $1$, $\mathrm{Pic}^{1}_{X/S}$ et le
morphisme $x \mapsto \mathcal{O}(x)$, défini par la diagonale de
$X \times_S X$, forment une solution du problème d'Albanese ; d'où un
isomorphisme
\[
c : B \xrightarrow{\ \sim\ } A = \mathrm{Pic}^{0}_{B/S},
\qquad c\bigl(\mathrm{cl}(y) - \mathrm{cl}(x)\bigr) = N_y N_x^{-1},
\]
l'\emph{autodualité} de la jacobienne. La page se demande si $c$ est
l'isomorphisme $\phi_{\Theta}$ de la polarisation thêta, et répond, après deux
ratures : « Non, c'est sans doute la \emph{symétrique} de la
$\Theta$-polarisation. » C'est exact : $c = -\phi_{\Theta}$.\footnote{C'est
un résultat classique, que l'on trouve par exemple chez Milne (« Jacobian
varieties », 1986) ; il dépend de la normalisation du faisceau de Poincaré, et
c'est celle des pages. La page ramène la question au cas d'un corps
algébriquement clos, écrit $u(t) = $ classe de $t.\Theta - \Theta$, et
commence le calcul pour $t = \mathrm{cl}((y) - (x))$ ; la page 24 ne porte que
ces deux lignes, suivies d'une formule biffée, et le reste du feuillet est
blanc. La réponse n'est donc pas démontrée dans le dossier.} C'est ce signe que
la page 8 annonçait par son « attention au signe ! ».

\pagerange{25}{27}
\section*{Schémas abéliens principalement polarisés (pages 25 à 27)}

Soit $A/S$ un schéma abélien, $\delta$ une section de $\mathrm{NS}_{A/S}$,
$P^{\delta} = \mathrm{Pic}^{\delta}_{A/S}$, torseur sous le dual
$B = \mathrm{Pic}^{0}_{A/S}$, et $\mathcal{L}$ le faisceau universel sur
$A \times_S P^{\delta}$ ; $A$ agit sur $P^{\delta}$ à travers $\tilde{\delta}$,
et c'est l'action par transport de structure des translations de $A$. Si
$\delta$ est une polarisation, les $H^{i}(\mathcal{L}_x)$ s'annulent pour
$i > 0$, et $\mathrm{pr}_{2*}\mathcal{L}$ est localement libre de rang
$\chi(\delta)$. Si $\chi(\delta) = 1$, c'est un module inversible, et la
section tautologique définit un diviseur canonique $D_{\delta}$ sur
$A \times_S P^{\delta}$, relatif sur $P^{\delta}$. Alors $\tilde{\delta}$ est
un isomorphisme, $P^{\delta}$ devient un $A$-torseur, $D_{\delta}$ est stable
par l'action diagonale de $A$ et descend au quotient
$(A \times_S P^{\delta})/A \simeq P^{\delta}$ en un diviseur $\Theta$, avec
$\chi(\Theta) = 1$.\footnote{« donc d'un » et « a la vertu » sont lus sans
certitude ; la page note $\tilde{\delta}$ puis $\bar{\delta}$ le même
homomorphisme.}

\textbf{Proposition} (pages 25 et 26). \emph{Soit $A/S$ un schéma abélien de
dimension relative $g$, et $\delta$ une polarisation de degré $1$.}
\begin{itemize}
\item \emph{a)} \emph{Il existe un $A$-torseur $P$ et un diviseur relatif positif
  $\Theta$ sur $P$, avec $\deg \Theta^{g}/g! = 1$, tels que $\delta$ soit la
  polarisation définie par $\Theta$. Le couple $(P, \Theta)$ est unique à
  isomorphisme unique près : une section de $A$ dont la translation préserve
  $\Theta$ est la section unité.}
\item \emph{b)} \emph{Si $P$ est un $A$-torseur et $\Theta$ un diviseur relatif
  positif avec $\deg \Theta^{g}/g! = 1$, alors $\Theta$ est ample relativement
  à $S$, et la polarisation qu'il définit sur $A$ est de degré $1$.}
\item \emph{c)} \emph{Sous b), soit $m : A \times_S P \to P$ l'action et
  $\Theta' = m^{-1}(\Theta)$. Alors $\Theta'$ est un diviseur relatif sur
  $A \times_S P$ au-dessus de $P$, et le morphisme $P \to \mathrm{Pic}_{A/S}$
  qu'il définit est un isomorphisme de $P$ sur $\mathrm{Pic}^{\delta}_{A/S}$,
  $\delta$ étant la polarisation définie par $\Theta$.}
\item \emph{d)} \emph{Si $A = \mathrm{Alb}^{0}_{X/S}$ pour une courbe $X/S$ propre et
  lisse de genre $g$, et $\delta$ la polarisation canonique, définie par le
  diviseur diagonal de $X \times_S X$, alors $\delta$ est de degré $1$ et
  $P^{\delta}$ est canoniquement isomorphe, comme $A$-torseur, à
  $\mathrm{Pic}^{g-1}_{X/S} = \mathrm{Alb}^{g-1}_{X/S}$, par l'unique
  isomorphisme qui envoie $\Theta$ sur $W_{g-1}$.}
\end{itemize}
\footnote{La première partie de la proposition ne porte pas de lettre sur la
page ; on la note a). Les énoncés sont ceux de la page ; la page ne donne pas
de preuve. Ils sont justes : en b), un diviseur effectif sur une variété
abélienne définit un fibré ample dès que $\chi \neq 0$, et
$\deg\Theta^{g}/g! = \chi$ par Riemann--Roch ; l'unicité de a) vient de ce que
le noyau de $\phi_{\Theta}$ est nul. En c), la page écrit « Alors $\Theta$ »,
là où il s'agit de $\Theta'$. En d), $W_{g-1}$ est l'image de
$X^{(g-1)}$, diviseur thêta de $\mathrm{Pic}^{g-1}$ par le théorème de Riemann ;
« comme $A$-torseur » et « ce qui le caractérise » sont lus sans certitude, et
une note encadrée de la marge dit « question de signe ! ».}

Un passage encadré et barré reprend ensuite le cadre général, avec une
correspondance divisorielle symétrique $C$ sur $X \times_S X$ définissant une
polarisation $\gamma$ de $A$ ; il est abandonné.

\emph{(c bis).} Soit $(P, \Theta)$ comme en b), et $P^{(-1)}$ le torseur
inverse ; l'identité des schémas sous-jacents est un isomorphisme
$\sigma_P : P \to P^{(-1)}$ compatible avec la symétrie de $A$. Posons
$\Theta^{-} = \sigma_P(\Theta)$. Comme $[-1]^{*}$ agit trivialement sur
$\mathrm{NS}$, $\Theta^{-}$ définit sur $A$ la même polarisation que $\Theta$,
donc, d'après a), il existe un unique isomorphisme de $A$-torseurs
\[
K : P^{(-1)} \xrightarrow{\ \sim\ } P, \qquad K(\Theta^{-}) = \Theta,
\]
qu'on peut voir comme une section $K \in \Gamma(S, P^{(2)})$.\footnote{La page
écrit $\Theta^{-*}$, exposant lu douteusement, et la flèche
$K : P \leftarrow P^{-1}$ avec $K(\Theta) = \Theta^{-*}$ ; on oriente $K$ et
la condition dans le même sens. La lettre de « (c bis) » peut aussi se lire
e.} Dans le cas de d), $P = \mathrm{Pic}^{g-1}_{X/S}$ et
$P^{(2)} = \mathrm{Pic}^{2g-2}_{X/S}$, et $K$ est la classe du fibré canonique
$\Omega^{1}_{X/S}$ : l'involution $L \mapsto \Omega^{1} \otimes L^{-1}$ de
$\mathrm{Pic}^{g-1}$ préserve $W_{g-1}$, par Riemann--Roch et la dualité de
Serre, et elle est équivariante de $P^{(-1)}$ vers $P$.\footnote{La
justification est nôtre ; la page affirme le résultat. Ses points fixes sont
les thêta-caractéristiques, dont le dossier ne parle pas.}

C'est la section de $\mathrm{Pic}^{2g-2}$ que la page 8 attendait (« la
section canonique $K$, à prouver ! »), obtenue cette fois sans le déterminant
de la cohomologie, par la seule symétrie du diviseur thêta. Une note de la
marge demande enfin de vérifier que c'est « la section canonique » de
$(\mathrm{Pic}^{\delta}_{A/S})^{(2)}$ pour toute section $\delta$ de
$\mathrm{NS}_{A/S}$ — c'est-à-dire de la comparer aux trivialisations de
carrés des pages 12 et 13, dont c'est le cas $\chi(\delta) = 1$. Le dossier
s'arrête sur cette demande.\footnote{L'exposant de $\mathrm{Pic}$ dans la note
de marge est illisible ; la lecture « $(\mathrm{Pic}^{\delta}_{A/S})^{(2)}$ »
est nôtre, d'après les pages 12 et 13.}

\end{document}
